\DocumentMetadata{
  pdfversion=2.0,
  pdfstandard=ua-2,
  lang=en-US,
  tagging=on,
  tagging-setup={math/setup=mathml-SE,tabsorder=structure}
}
\documentclass[11pt,letterpaper]{article}
\usepackage[margin=0.68in,includefoot]{geometry}
\usepackage{newcomputermodern}
\usepackage{amsmath,amscd,mathtools}
\usepackage{tikz,adjustbox}
\providecommand{\square}{\mdlgwhtsquare}
\usepackage[hidelinks]{hyperref}
\hypersetup{
  pdftitle={Topology PhD Exam, September 2000},
  pdfauthor={University of Florida Department of Mathematics},
  pdfsubject={Graduate mathematics examination},
  pdfdisplaydoctitle=true,
  bookmarksopen=true
}
\setcounter{secnumdepth}{0}
\setlength{\parindent}{0pt}
\setlength{\parskip}{0.28em}
\setlength{\emergencystretch}{3em}
\setlength{\leftmargini}{2.15em}
\setlength{\leftmarginii}{2.65em}
\setlength{\leftmarginiii}{3em}
\renewcommand{\labelenumi}{\textbf{\theenumi.}}
\pagestyle{empty}
\raggedbottom
\allowdisplaybreaks
\newenvironment{examproblems}[1]{%
  \begin{enumerate}%
  \setcounter{enumi}{#1}%
  \setlength{\topsep}{0.35em}%
  \setlength{\partopsep}{0pt}%
  \setlength{\itemsep}{0.48em}%
  \setlength{\parsep}{0.18em}%
}{\end{enumerate}}
\newenvironment{examsubparts}{%
  \begin{enumerate}%
  \setlength{\topsep}{0.18em}%
  \setlength{\partopsep}{0pt}%
  \setlength{\itemsep}{0.16em}%
  \setlength{\parsep}{0.1em}%
}{\end{enumerate}}
\newenvironment{examinstructions}{%
  \par\begingroup\itshape
}{%
  \par\endgroup\vspace{0.18em}
}
\makeatletter
\renewcommand\section{\@startsection{section}{1}{0pt}%
  {-1.25ex plus -.25ex minus -.1ex}%
  {0.55ex plus .1ex}%
  {\normalfont\large\bfseries}}
\renewcommand{\@maketitle}{%
  \newpage\null\vskip -1em%
  {\footnotesize\bfseries
    \href{https://www.ufl.edu/}{University of Florida}, \href{https://math.ufl.edu/}{Department of Mathematics}\par}%
  \vskip -0.5em%
  \tagstructbegin{tag=Title}%
    {\LARGE\bfseries\href{https://gma.math.ufl.edu/exams/topology-phd/}{Topology PhD Exam}, September 2000\par}%
  \tagstructend%
  \vskip 0.25em%
  \tagmcbegin{artifact}%
    \hrule height 1pt%
  \tagmcend%
  \vskip 1em%
}
\makeatother
\title{Topology PhD Exam, September 2000}
\author{}
\date{}
\begin{document}
\maketitle
\thispagestyle{empty}
\begin{examinstructions}
Work the following problems and show all work. Support all statements to the best of your ability.
\end{examinstructions}
\begin{examproblems}{0}
\item
State and prove the Brouwer Fixed Point Theorem. You may assume some homology.
\item
State and prove the Baire Category Theorem.
\item
Give the Mayer–Vietoris sequence for homology. Use this to compute the homology groups for the~\(n\)-sphere \(S^n\).
\item
Suppose that \(f,g:X\to P\) are two maps and that~\(P\) is a compact connected polyhedron. Show that there is an \(\epsilon>0\) such that if \(d(f(x),g(x))<\epsilon\) for all \(x\in X\), then~\(f\) and~\(g\) are homotopic.
\item
Show that \(\pi_1(S^1)\cong Z\), where~\(Z\) is the group of integers and \(S^1\) is the circle.
\end{examproblems}
\begin{examinstructions}
Answer the following with complete definitions or statements or short proofs.
\end{examinstructions}
\begin{examproblems}{5}
\item
Define retraction and deformation retraction.
\item
State Urysohn’s Lemma. State the Tietze Extension Theorem.
\item
State the Five Lemma.
\item
Define the cone of a space~\(X\). Define the suspension of~\(X\).
\item
State the exact homology sequence for a pair of spaces.
\item
State the Jordan Curve Theorem.
\item
Let~\(A\) be a countable subset of \(R^2\). Show that \(R^2\setminus A\) is arcwise connected.
\item
Show that the Cantor set less its endpoints is homeomorphic to the irrationals in the unit interval. [Hint: Use the Cantor Ternary Function.]
\item
Let~\(V\) be a complete metric space such that no point is isolated. Can~\(V\) be a countably infinite set?
\item
Let \(0\to G\xrightarrow{g}H\xrightarrow{h}J\to 0\) be an exact sequence. What does this say about~\(g\) and~\(h\)? Prove your claims.
\end{examproblems}
\end{document}
